% REVIEW DRAFT: not an input of the current published manuscript.
% This file is not proof certification; see geometry-audit-2026-11-21.json.

\section{From a zero line constant term to cuspidality}
\label{sec:sr-theta-line-cuspidality}

Here is a local-to-global consequence of the actual split matching
representation. It replaces a separate assumption that a zero line
constant term already implies cuspidality.

Let \(W\) be the fixed target of rank \(d=a+a'\), and suppose
\[
                    \delta=(a-a')/2>3.
\]
At a good fully split finite place \(v\), put \(k=F_v\) and
\(q_v=|\mathcal O_k/\mathfrak p_k|\).
In the specified split action the common target representation is,
up to its already specified unitary determinant twist,
\[
 \pi_v=
 i_{(a,a')}^{\,\mathrm{GL}_d(k)}
       \bigl((\xi\circ\det_a)\boxtimes
                         (\xi'\circ\det_{a'})\bigr),
 \tag{U1}\label{eq:sr-cusp-line-local-representation}
\]
where \(i\) is normalized induction and \(\xi,\xi'\) are
unramified unitary characters. The order of these two characters
will not matter for the argument. This is an irreducible
representation by the local calculation already proved.

\begin{srlemma}[Generic pieces of the split Jacquet module]
\label{lem:sr-cusp-line-jacquet}
For \(1\le r\le d/2\), take the normalized Jacquet module of
\(\pi_v\) for the three-block parabolic
\((r,d-2r,r)\).
Consider its first \(\mathrm{GL}_r(k)\) factor.
It has no generic irreducible constituent if \(r>2\).
If \(r=2\), every generic constituent on that factor is
\[
 i_{(1,1)}^{\,\mathrm{GL}_2(k)}
   \left(\xi|\cdot|^{-(a-1)/2}\boxtimes
                      \xi'|\cdot|^{-(a'-1)/2}\right),
 \tag{U2}\label{eq:sr-cusp-line-GL2-piece}
\]
up to permutation and a common determinant twist.
In particular the absolute values of the two inducing characters
have ratio \(q_v^\delta\) at a uniformizer, or its inverse.
\end{srlemma}

\begin{proof}
We spell out the relevant geometric-lemma calculation.
The double cosets are indexed by matrices
\((m_{ij})\) of nonnegative integers, with row sums
\((a,a')\) and column sums \((r,d-2r,r)\).
The corresponding subquotient is the tensor product, by columns,
of normalized inductions of the normalized Jacquet modules
of the two determinant characters.

In the first column its two block sizes are \(m_{11},m_{21}\)
with sum \(r\). A determinant representation of
\(\mathrm{GL}_m(k)\) is nongeneric for \(m\ge2\):
its upper unipotent action is trivial whereas a nondegenerate
Whittaker character is nontrivial. Its Whittaker coinvariant
space is zero.
Whittaker coinvariants of normalized induction are the tensor
product of those of the inducing blocks. One may verify this
by the Bruhat filtration: only the open cell survives the
nondegenerate character; integrating along that cell gives the
tensor product. Twisted unipotent coinvariants are exact over
\(\mathbb C\), by averaging on compact open subgroups and
passing to their increasing union. Therefore vanishing on
these subquotients proves vanishing for the entire Jacquet
module and for every one of its subquotients.
A generic first-column constituent thus requires
\(m_{11},m_{21}\le1\), so \(r\le2\).

If \(r=2\), both entries must be \(1\).
In a normalized Jacquet module the determinant character of
\(\mathrm{GL}_a\), restricted to its first block of size one,
is \(\xi|\cdot|^{-(a-1)/2}\): the ordinary action is \(\xi\),
and normalization multiplies by the inverse square root of
the parabolic module. This power is independent of the
distribution of the remaining \(a-1\) coordinates into the
later two columns. The corresponding primed power is
\(-(a'-1)/2\), proving (U2).
Its character ratio has norm exponent \(\delta\).
This principal series is irreducible: its ratio cannot equal
\(|\cdot|^{\pm1}\), since \(\delta>3\).
Thus there is no additional generic constituent with different
parameters hidden in (U2).

This is precisely the ordinary Bruhat geometric lemma of
Bernstein--Zelevinsky, \emph{Induced representations of reductive
\(p\)-adic groups I}, \S2 (Lemma~2.12 and its filtration);
the block computation just given identifies all the terms
needed here.
\footnote{\url{https://www.math.tau.ac.il/~bernstei/Publication_list/publication_texts/BZ-induced-1.pdf}.}
\end{proof}

\begin{srlemma}[The two-dimensional piece cannot be essentially unitary]
\label{lem:sr-cusp-line-unitarity}
An irreducible representation of the form (U2) is not essentially
unitary when \(\delta>3\).
\end{srlemma}

\begin{proof}
A common unramified real determinant twist making its central
character unitary makes the two normalized Satake parameters
have absolute values \(q_v^{\delta/2}\) and
\(q_v^{-\delta/2}\).
If the representation were unitary after such a twist, normalize
the spherical vector to have norm one.
The Hecke operator of
\(K\operatorname{diag}(\varpi,1)K\), with \(\operatorname{vol}K=1\),
is a sum of \(q_v+1\) unitary right translations.
Its norm is at most \(q_v+1\).
Directly decomposing this double coset into its \(q_v+1\)
Iwasawa cosets gives its spherical eigenvalue
\[
                 q_v^{1/2}(\lambda_1+\lambda_2).
\]
Therefore
\[
 |\lambda_1+\lambda_2|
          \le q_v^{1/2}+q_v^{-1/2}.                 \tag{U3}
 \label{eq:sr-cusp-line-unitary-Hecke}
\]
But the reverse triangle inequality gives
\[
 |\lambda_1+\lambda_2|
       \ge q_v^{\delta/2}-q_v^{-\delta/2}
       > q_v^{3/2}-q_v^{-3/2}
       >q_v^{1/2}+q_v^{-1/2}.
\]
For the last strict inequality, multiply by \(q_v^{3/2}\)
and use
\[
 q_v^3-q_v^2-q_v-1
       =(q_v-2)(q_v^2+q_v+1)+1>0
       \qquad(q_v\ge2).
\]
This contradicts (U3).
A unitary character twist changes neither the norm comparison
nor this conclusion. Thus no determinant twist can unitarize
(U2).
\end{proof}

\begin{srproposition}[The actual line-to-cusp implication]
\label{prop:sr-theta-line-implies-cusp}
Let \(F\) be a smooth moderate automorphic form in the span of
the actual large and small polynomial theta-period images
for a fixed triple. Suppose its constant term along the
unipotent radical of an isotropic line is identically zero.
Then \(F\) is cuspidal.

More generally the assertion holds for such a form whose local
module at one good fully split place is a quotient of a finite
sum of copies of (U1).
\end{srproposition}

\begin{proof}
All rational isotropic lines are conjugate by \(U(W)(F)\),
by Witt extension. Automorphy therefore gives the zero constant
term for every rational isotropic line. The zero assertion
also holds for all right translates of \(F\).

If \(F\) were not cuspidal, choose the least \(r\ge1\) for
which its constant term along a rational isotropic \(r\)-plane
is nonzero. Necessarily \(r\ge2\).
Its Levi is
\(\mathrm{GL}_r(\mathbb A_E)\times U(W_0)(\mathbb A_F)\).
Every proper general-linear constant term of that nonzero
Levi function is zero. Indeed staged unipotent integration
identifies it with the constant term along a flag
\(L\subset L_r\), \(0<\dim L<r\); integrate first over
the radical for \(L\). That first integral is zero by
minimality of \(r\). All the unipotent quotients here are
compact, so the staged integrals are ordinary Fubini integrals.
Thus a nonzero general-linear slice of this Levi function
is a cuspidal automorphic form on \(\mathrm{GL}_r(\mathbb A_E)\).

We recall two classical general-linear facts and their scope.
Cuspidal automorphic representations of
\(\mathrm{GL}_r(\mathbb A_E)\) are globally generic and are
essentially unitary. Genericity follows by successive
Fourier expansions along the last-column unipotent groups:
cuspidality removes the zero coefficient, and the rational
general linear group is transitive on their nonzero
characters. Iterating gives a nonzero nondegenerate Whittaker
coefficient. This is the usual cuspidal Whittaker expansion
of Shalika.
\footnote{J.~Shalika, \emph{The multiplicity one theorem for
\(\mathrm{GL}_n\)}, Ann.~of Math.~100 (1974), 171--193:
\url{https://annals.math.princeton.edu/1974/100-2/p01}.}
After fixing a central character, its absolute value is a real
power of the idele norm; twisting by the corresponding real
power of determinant makes the central character unitary.
Cuspidal rapid decrease modulo center then realizes the
resulting representation unitarily in the cuspidal \(L^2\)
space. In particular each local factor is essentially unitary.

These facts apply to the nonzero slice just obtained.
For precision one may take a nonzero irreducible constituent
of its cuspidal automorphic span. The theta images are compact
finite and infinitesimal-character finite at infinity by
their actual real matching quotients. Restriction of a
finite-dimensional upper compact type to the Levi compact
group is still compact finite. To check the center condition
on a constant term, use PBW projection along the parabolic
nilpotent algebra. On a function constant along that radical,
the upper center acts by its Harish--Chandra projection to
the Levi center, with the ordinary half-modulus shift.
Under the Harish--Chandra isomorphisms this is the inclusion
\(\operatorname{Sym}(\mathfrak h)^{W_G}
  \subset\operatorname{Sym}(\mathfrak h)^{W_M}\).
The second algebra is finite over the first, because the
Weyl groups are finite. Thus a finite-codimensional
upper infinitesimal-character annihilator gives a
finite-codimensional annihilator for the Levi center.
In particular the general-linear slice has only finitely
many generalized infinitesimal central weights.

One can now take an actual central eigencomponent:
decompose the norm-one idele class center, which is compact,
into characters, and take a nonzero component. Its compact
Fourier weights and its remaining real exponential weights
are finite by the just-proved center condition. If a
nilpotent logarithmic factor occurs in the real central
direction, apply its highest nonzero central derivative
to obtain a nonzero eigenvector. These operations commute
with the general-linear group and preserve both cuspidality
and the local Jacquet-module restriction.
After the determinant norm twist making this character
unitary, the slice lies in the cuspidal \(L^2\) space modulo
center. Its usual discrete cuspidal decomposition, or the
finite-dimensional fixed-level, fixed-compact-type,
fixed-infinitesimal-character cusp space, gives a nonzero
irreducible cuspidal constituent. Its Whittaker expansion
gives genericity. This passage uses the general-linear
cuspidal spectral theory only; it does not assume a
multiplicity-one theorem for the unitary group.

At the chosen split place, a parabolic stabilizing an
\(r\)-plane becomes the three-block parabolic
\((r,d-2r,r)\) in \(\mathrm{GL}_d(k)\).
The two \(\mathrm{GL}_r(k)\) factors of its Levi are the
two completions of \(\mathrm{GL}_r(E_v)\).
The nonzero cuspidal constituent must be generic and
essentially unitary at each of them.
On the other hand its first factor is a subquotient of
the Jacquet module calculated in the preceding lemma.
Using an unnormalized constant term instead of normalized
Jacquet restriction contributes only a common determinant
norm twist; it does not affect genericity or essential
unitarity.

For \(r>2\) this contradicts the absence of a generic
constituent. For \(r=2\) it contradicts the nonunitarity
lemma. Hence no such \(r\) exists.
Every proper rational parabolic is contained in one
stabilizing a nonzero isotropic subspace, and staged
integration then gives zero for all of its constant terms.
This is cuspidality.

For an actual triple the local period factorization gives
exactly the representation (U1) at each good fully split
place; for a sum or difference there are at most finitely
many copies. Such places exist by the good-prime argument.
Finally the exact rank inequality
\(2(p+q)+1<(a-a')/2\) gives \(\delta>3\)
whenever the target has an isotropic line, since then
\(p+q\ge1\). No integrality of \(\delta\) is used. When it has
none, its quotient has no proper rational parabolic and
the assertion is immediate.
\end{proof}

This proposition supplies the line-to-cusp step only. In the
signature induction one must still show that the other Levi
component of the difference vanishes when the first chosen
component has been matched. The two-step contraction and
lower-rank injectivity address that separate statement.
